% Source: source/普通高考/2012/2012山东理.pdf, page 1, question 6.
% Nodes: start; input a; initialize P=0,Q=1,n=0; decision P<=Q?;
% yes branch P=P+a^n -> Q=2Q+1 -> n=n+1; no branch output n -> finish.
% Directed edges: start->input->initialize->decision; decision(是)->P update
% ->Q update->n update->return to the decision stem; decision(否)->output->finish.
% The return path is the independent outer left loop and joins the main stem
% immediately above the decision; no other branches share a segment.
% All edges are solid black with arrowheads; labels 是/否 sit beside their branch.

\begin{tikzpicture}[
  >=Stealth,
  x=1cm,
  y=0.84cm,
  line width=0.45pt,
  line join=round,
  line cap=round,
  font=\normalsize,
  flowtext/.style={anchor=center, align=center,
    execute at begin node={\setlength{\parindent}{0pt}}},
  every node/.style={flowtext},
  flowbox/.style={flowtext, draw, minimum height=0.68cm,
    inner xsep=4pt, inner ysep=2.5pt},
  branchlabel/.style={flowtext, inner sep=0pt},
  terminal/.style={flowbox, rounded corners=0.18cm, minimum width=1.35cm},
  process/.style={flowbox},
  decision/.style={flowbox, diamond, aspect=3.0, minimum width=3.55cm,
    minimum height=0.95cm},
  io/.style={flowbox, trapezium, trapezium left angle=72,
    trapezium right angle=108, minimum width=2.25cm, minimum height=0.76cm},
]
  \path[use as bounding box] (-2.35,0.25) rectangle (5.65,10.65);

  \node[terminal] (start) at (0,10.00) {开始};
  \node[io] (input) at (0,8.55) {输入 $a$};
  \node[process, minimum width=4.35cm] (initialize) at (0,7.05)
    {$P=0,Q=1,n=0$};
  \node[decision] (test) at (0,5.20) {$P\leqslant Q$};
  \node[process, minimum width=3.35cm] (updateP) at (0,3.62)
    {$P=P+a^n$};
  \node[process, minimum width=3.35cm] (updateQ) at (0,2.55)
    {$Q=2Q+1$};
  \node[process, minimum width=3.35cm] (updateN) at (0,1.12)
    {$n=n+1$};
  \node[io] (output) at (4.00,3.62) {输出 $n$};
  \node[terminal] (finish) at (4.00,2.55) {结束};

  % Main flow and the single downward arrow entering the decision.
  \draw[->] (start.south) -- (input.north);
  \draw[->] (input.south) -- (initialize.north);
  \coordinate (merge) at (0,6.15);
  \draw (initialize.south) -- (merge);
  \draw[->] (merge) -- (test.north);

  % Yes branch and its sequential updates.
  \draw[->] (test.south) --
    node[branchlabel, anchor=west, xshift=3pt] {是} (updateP.north);
  \draw[->] (updateP.south) -- (updateQ.north);
  \draw[->] (updateQ.south) -- (updateN.north);

  % Independent loop-back path to the decision stem; its arrow ends at the
  % shared merge before the independent downward stem arrow.
  \draw (updateN.south) -- (updateN.south |- 0,0.55)
    -- (-2.05,0.55);
  \draw (-2.05,0.55) -- (-2.05,6.15);
  \draw[->] (-2.05,6.15) -- (-0.55,6.15) -- (merge);

  % No branch exits right and terminates independently.
  \draw[->] (test.east) -- node[branchlabel, above] {否}
    (4.00,5.20) -- (4.00,4.55) -- (output.north);
  \draw[->] (output.south) -- (finish.north);
\end{tikzpicture}
